\chapter*{Abstract}

This thesis addresses the critical vulnerability of deep learning models to adversarial perturbations, particularly the limitation where standard defense mechanisms fail to generalize across different threat models. Our research proposes a novel perspective that treats distinct families of adversarial attacks as separate "domains," aiming to identify and quantify the "cross-attack robustness gap" inherent in current training paradigms.

\noindent We establish a comprehensive experimental framework to evaluate the transferability of robustness across diverse adversarial environments. By subjecting models trained on specific perturbation norms to a wide spectrum of unseen attack algorithms, the study systematically measures the performance degradation that occurs when shifting between these domains.

\noindent Ultimately, this quantitative analysis serves as the empirical foundation for the subsequent phase of research. By characterizing the domain shifts between attack families, we pave the way for integrating Domain Generalization principles into adversarial training. This progression aims to directly mitigate the observed robustness gaps, ensuring that future models can effectively generalize to unforeseen and heterogeneous adversarial scenarios.